%% Appendix B -- from augur_AMT_부록B draft v1 (2026-09-22). Synthetic / algebraic only: no \fieldnum.
%% Phase 3: typeset as Appendix B, between Appendix A (configuration) and Appendix C (calibration-block length).
\section{The projected Jacobian, column scaling and the joint covariance}
\label{app:B}

\subsection{Setup}
\label{app:B1}

Write the forward model of Sect.~\ref{sec:2.1} for one spectrum as
\begin{equation*}
\alpha \;=\; \Phi(\theta)\,c \;+\; \varepsilon ,
\end{equation*}
where $\theta = (\text{shift}, \text{squeeze})$ enters the columns of $\Phi$
through the resampling of the reference cross sections, and $c$ collects the
gas amounts and the Chebyshev baseline coefficients. The problem is separable:
for any $\theta$ the linear parameters are determined by least squares,
$c^\ast(\theta) = \Phi^{+}\alpha$, and the objective reduces to the \emph{profiled}
residual
\begin{equation*}
r(\theta) \;=\; \bigl(I - \Phi(\theta)\Phi(\theta)^{+}\bigr)\,\alpha
\;=\; P^{\perp}(\theta)\,\alpha .
\end{equation*}
The optimiser searches only $\theta$; $c$ is never a search direction. This is
variable projection \citep{Golub_1973}.

\subsection{The derivative has two terms}
\label{app:B2}

Differentiating the projector rather than the fit gives, for each nonlinear
parameter $\theta_j$ and with $\Phi_j = \partial\Phi/\partial\theta_j$,
\begin{equation*}
\frac{\partial r}{\partial \theta_j}
\;=\; -\Bigl[\, P^{\perp}\Phi_j\Phi^{+}
\;+\; \bigl(P^{\perp}\Phi_j\Phi^{+}\bigr)^{\!\top} \Bigr]\alpha .
\end{equation*}
The first term is the change in the residual caused by moving the column
space; the second is the change caused by the reprojection of $\alpha$ onto
it. \citet{Kaufman_1975} drops the second, which is cheaper and is often adequate
because it does not change the stationary point -- both expressions vanish
where the true gradient vanishes -- only the path taken to it and the
curvature reported there.

It is not adequate here. On a small separable problem with the same structure
as Sect.~\ref{sec:2.1} -- two Gaussian absorber columns and three Chebyshev columns -- the
two-term expression agrees with a central finite difference of $r(\theta)$ to
$2\times10^{-10}$ relative, while the one-term expression is in error by 20\,\%
and 17\,\% in the two nonlinear parameters. The step-size sweep of Sect.~\ref{sec:3.2} measures
the same thing on the real problem: the analytic derivative descends to the
round-off floor, which an expression missing a term of that size could not do.

\subsection{Column scaling changes nothing, until it changes everything}
\label{app:B3}

Replace $\Phi$ by $\Phi D$ for any invertible diagonal $D$. Then
$(\Phi D)^{+} = D^{-1}\Phi^{+}$, so $P = \Phi D D^{-1}\Phi^{+} = \Phi\Phi^{+}$
and $P^{\perp}$ are unchanged, and in the derivative of Appendix~\ref{app:B2} the factors
$D$ and $D^{-1}$ cancel term by term. The profiled residual, its derivative
and the fitted $\theta$ are therefore invariant under rescaling of the
columns; only $c$ and its covariance rescale, and they rescale exactly
inversely, so the retrieved amounts are unchanged.

That is the algebra. The arithmetic is not invariant, and the failure is
abrupt rather than gradual, because $\Phi^{+}$ is computed with a rank
tolerance (Table~\ref{tab:colscaling}).

\begin{table}[t]
\caption{Effect of rescaling one column of the small test problem of Appendix~\ref{app:B2} on the condition number $\kappa$, on the projector $P^{\perp}$ and on the derivative $\partial r/\partial\theta$.}
\label{tab:colscaling}
\begin{tabular}{llll}
\tophline
column scaling & $\kappa(\Phi D)$ & $\lVert\Delta P^{\perp}\rVert$ & relative change in $\partial r/\partial\theta$ \\
\middlehline
none & $9.3$ & $0$ & $0$ \\
$10^{-3}$ on one column & $6.1\times10^{5}$ & $2.4\times10^{-10}$ & $1.1\times10^{-10}$ \\
$10^{-8}$ & $6.1\times10^{11}$ & $1.3\times10^{-4}$ & $5.1\times10^{-5}$ \\
$10^{-12}$ & $7.7\times10^{15}$ & $1.00$ & $0.48$ \\
$10^{-19}$ & $2.1\times10^{16}$ & $1.00$ & $0.48$ \\
\bottomhline
\end{tabular}
\end{table}

Once the scaling drives the condition number past the tolerance the
pseudo-inverse discards the column, the projector is onto a different
subspace, and the invariance is lost completely rather than degraded. The
operational problem sits at the bottom of that table by construction: a
cross-section column is of order $10^{-19}$\,cm$^2$ and a Chebyshev column is
of order unity. Column normalisation is therefore not a numerical nicety but
the condition under which the invariance of Appendix~\ref{app:B3} holds at all, and the
retrieval normalises every column to unit norm before the solve.

\subsection{The joint covariance and what the fit usually reports}
\label{app:B4}

Linearising the full problem in $(c,\theta)$ at the solution gives the
Jacobian $J = [\,\Phi \;\; B\,]$ with
$B = \partial(\Phi c)/\partial\theta$, and the covariance
$\sigma^{2}(J^{\top}J)^{-1}$. Partitioning it, the $\theta$ block is the
Schur complement
\begin{equation*}
\mathrm{Cov}(\theta) \;=\; \sigma^{2}\bigl(B^{\top}P^{\perp}B\bigr)^{-1},
\end{equation*}
which is the projected curvature of Appendix~\ref{app:B2} and not $\sigma^2(B^\top B)^{-1}$: the
part of $B$ that lies in the column space of $\Phi$ is absorbed by the linear
parameters and constrains $\theta$ not at all. The $c$ block is correspondingly
larger than the $(\Phi^{\top}\Phi)^{-1}$ that a fit with $\theta$ held fixed
would return, by an amount that depends on how much of each column's
sensitivity the nonlinear parameters can mimic. On the toy problem above the
inflation runs from 1.001 to 1.257 across the five columns.

Both statements are elementary and neither is new. What matters for Sect.~\ref{sec:5} is that
a retrieval which reports $\sigma^2(\Phi^{\top}\Phi)^{-1}$ -- the covariance of
the linear parameters at a fixed $\theta$ -- is reporting a quantity that is
correct only if $\theta$ was known rather than fitted, and the difference is
not small when $\theta$ is poorly determined. Augur forms the full
$(J^{\top}J)^{-1}$ and reports both blocks; the $\theta$ block is available at no extra cost once
the projected Jacobian of Appendix~\ref{app:B2} has been formed. It describes only the local curvature at the
solution, which is why the identifiability statistics of Sect.~\ref{sec:5} are taken from profile scans of the
projected objective rather than from this block.

\subsection{Implementation of the projected Jacobian and relation to DOAS codes}
\label{app:B5}

In the implementation of Sect.~\ref{sec:2.3} the weighted, column-scaled design matrix is factorised as
$\mathbf{A}\mathbf{S}^{-1}=\mathbf{Q}\mathbf{R}$, which delivers $\hat{\mathbf{c}}$, the residual, and the projector
$\mathbf{P}^{\perp}=\mathbf{I}-\mathbf{Q}\mathbf{Q}^{\!\top}$ from one factorisation, with
$(\mathbf{A}^{+})^{\!\top}=\mathbf{Q}\mathbf{R}^{-\top}$ in Eq.~(\ref{eq:6}). Only the gas columns of
$\mathbf{D}_k$ are non-zero: differentiating Eq.~(\ref{eq:3}),
\begin{equation}
\frac{\partial}{\partial \delta_i}\,\sigma_i(p_i') = \sigma_i'(p_i'),
\qquad
\frac{\partial}{\partial s_i}\,\sigma_i(p_i') = (p-p_c)\,\sigma_i'(p_i'),
\label{eq:7}
\end{equation}
so $\sigma_i'$ is obtained once from the derivative of the interpolating spline that
already represents each reference. The polynomial and etalon columns are independent of
$\theta$ and contribute nothing. Both terms of Eq.~(\ref{eq:6}) cost one triangular solve and
one multiplication by $\mathbf{Q}$; their derivation and the invariance under the column scaling $\mathbf{S}$ are
given in Appendices~\ref{app:B2} and \ref{app:B3}. The gradient of the objective is the same with and without
the second term, because $\mathbf{A}^{+}\mathbf{r} = 0$; the second term adds the exact curvature and makes the
Jacobian testable against finite differences to round-off (Sect.~\ref{sec:3.2}). We have not compared the
convergence rates of the two forms.

The separable structure is not new to DOAS; what the established treatments do not form is the derivative of
the linear solution with respect to the nonlinear parameters. QDOAS solves the slant columns and the polynomial
by singular value decomposition inside each evaluation of the fitting function while its
Marquardt--Levenberg loop iterates over the nonlinear parameters only -- the structure WinDOAS itself named
``Marquardt-SVD'' -- but its source states that the derivatives with respect to shift and stretch are ``always
numeric'', so no term of the form of Eq.~(\ref{eq:6}) is formed. DOASIS \citep{Kraus_2006} goes the other way:
once shift and squeeze are fitted, the concentrations are handed to the nonlinear solver as well, and the thesis
notes that the algorithm may then fail to find the global optimum that the linear step would have guaranteed.
The standard textbook scheme \citep[][Sect.~8.3.4]{Platt} alternates one linear solve with one
Levenberg--Marquardt step, explicitly ``with fixed parameters $a_j$ and $c_h$'' -- block-coordinate
optimisation, in which $\partial \hat{\mathbf{c}}/\partial\theta$ does not appear. Outside DOAS, the analytic form
of Eq.~(\ref{eq:6}) is established in atmospheric remote sensing: \citet{B_rligea_2023} use exactly this
projected derivative, with the Kaufman variant, for a SWIR retrieval in which -- as in BIRRA
\citep{Gimeno_Garc_a_2011} -- the separability is inverted relative to Eq.~(\ref{eq:2}).

For the covariance of Sect.~\ref{sec:2.5}: the inverse of the gas block of Eq.~(\ref{eq:9}) is the Schur complement
\begin{equation*}
\mathbf{A}_{\mathbf{W}}^{\top}\mathbf{A}_{\mathbf{W}} - \mathbf{A}_{\mathbf{W}}^{\top}\mathbf{V}(\mathbf{V}^{\top}\mathbf{V})^{-1}\mathbf{V}^{\top}\mathbf{A}_{\mathbf{W}},
\end{equation*}
which is never larger than the matrix inverted in Eq.~(\ref{eq:8}), so the joint standard errors are never
smaller than the conditional ones. $\mathbf{V}$ holds $\mathbf{c}$ fixed and is therefore not the projected
Jacobian of Eq.~(\ref{eq:6}); using Eq.~(\ref{eq:6}) there would count the linear block twice.

\subsection{Scaling of the fitted quantity and options not used operationally}
\label{app:B6}

Equation~(\ref{eq:4}) is solved on $\alpha$ scaled to order unity, and the linear coefficients
are scaled back afterwards. Ambient $\alpha$ in these windows is of order $10^{-7}\,\mathrm{cm^{-1}}$, so
residuals and gradients are of order $10^{-15}$, while the default termination tolerances of a general-purpose
least-squares routine are of order $10^{-8}$. On unscaled $\alpha$ the termination test can then be satisfied at
the starting point: the solver returns after a single function evaluation, having moved the shift by nothing,
and reports success. In one operational spectrum the unscaled fit exited with shift $+0.0000$\,px after
$n_{\mathrm{fev}}=1$, while the same spectrum scaled by $10^{7}$ converged to $-0.5000$\,px after
$n_{\mathrm{fev}}=5$. The failure is silent -- it produces a plausible number with a success flag -- and it is
specific to re-implementations in scaled, tolerance-driven optimisers rather than to the DOAS method itself.

The linear step can be solved with a non-negativity constraint on the gas coefficients, and a Tikhonov penalty
can be applied to the gas, etalon and auxiliary columns (never to the Chebyshev background, which would distort
it). Neither is used in the configuration reported here: the operational $\lambda$ is zero and the concentrations
are left unconstrained so that values near zero remain unbiased. The analytic Jacobian of Eq.~(\ref{eq:6})
requires the unconstrained linear step: at an active non-negativity bound the projector is not differentiable,
and the implementation falls back to finite differences in that case. Robust (Tukey bisquare) iteratively
reweighted fitting is available and likewise unused here.
